\documentclass[pdflatex,compress,9pt,
	xcolor={dvipsnames,dvipsnames,svgnames,x11names,table},
	hyperref={colorlinks = true,breaklinks = true, urlcolor = NavyBlue, breaklinks = true}]{beamer}
\usetheme{Goettingen}

\usepackage[utf8]{inputenc}
\usepackage[T2A,T1]{fontenc}
\usepackage{gensymb} % degree symbol
\usepackage[super]{nth}


% ----------------------------------------------------------------------------
% *** START BIBLIOGRAPHY <<<
% ----------------------------------------------------------------------------
\usepackage[
	backend=biber, 
	style = numeric,
%	style=nature,
%	style=science,
%	style=apa,
%	style=mla,
%	style=phys, 
	maxbibnames=99,
%	citestyle=authoryear,
	citestyle=numeric,
	giveninits=true,
	isbn=true,
	url=true,
	natbib=true,
	sorting=ndymdt,
	bibencoding=utf8,
	useprefix=false,
	language=auto, 
	autolang=other,
	backref=true,
	backrefstyle=none,
	indexing=cite,
]{biblatex}
\DeclareSortingTemplate{ndymdt}{
  \sort{
    \field{presort}
  }
  \sort[final]{
    \field{sortkey}
  }
  \sort{
    \field{sortname}
    \field{author}
    \field{editor}
    \field{translator}
    \field{sorttitle}
    \field{title}
  }
  \sort[direction=descending]{
    \field{sortyear}
    \field{year}
    \literal{9999}
  }
  \sort[direction=descending]{
    \field[padside=left,padwidth=2,padchar=0]{month}
    \literal{99}
  }
  \sort[direction=descending]{
    \field[padside=left,padwidth=2,padchar=0]{day}
    \literal{99}
  }
  \sort{
    \field{sorttitle}
  }
  \sort[direction=descending]{
    \field[padside=left,padwidth=4,padchar=0]{volume}
    \literal{9999}
  }
}

\addbibresource{YO.bib}
\renewcommand*{\bibfont}{\tiny} %\scriptsize \footnotesize

\setbeamertemplate{bibliography item}{\insertbiblabel}

% ----------------------------------------------------------------------------
% *** END BIBLIOGRAPHY <<<
% ----------------------------------------------------------------------------

\title{Economic assessment of landslide risk for the Waidhofen a.d. Ybbs region, Alpine Foreland, Lower Austria}
\author{Polina Lemenkova}
         \date{June 7, 2024}

\begin{document}
\begin{frame}
           \titlepage
\end{frame}

\section{Introduction}
\subsection{Summary}

\begin{frame}{National Civil Protection Management in Low Austria: ArcGIS-based approach for landslide assessment}
\begin{minipage}[0.4\textheight]{\textwidth}
\begin{columns}[T]
\begin{column}{0.5\textwidth}
\vspace{2em}
\begin{figure}[H]
	\centering
		\includegraphics[width=4.5cm]{Fig_00.jpg}
\end{figure}
\begin{figure}[H]
	\centering
		\includegraphics[width=2.5cm]{ArcGIS_logo.png}
\end{figure}
\end{column}
\begin{column}{0.60\textwidth}
\begin{enumerate}[i]
\vspace{1em}
\footnotesize
            \item \textcolor{red}{Research goals}: Improving the national landslide disaster control by evaluating the expected number of people affected by landslides.
            \item \textcolor{red}{Aim}: to assess potential risk and estimate economic damage caused by landslides. 
            \item \textcolor{red}{Study area}: Ybbs valley, Lower Austria
            \item \textcolor{red}{Software}: Arc GIS-based spatial analysis and estimation of the potential monetary losses caused by landslides.
            \item \textcolor{red}{Technical approach}: Geospatial and economic data processing by Arc GIS using spatial and economic analysis
            \item \textcolor{red}{Methods}: Defining elements at risk located in the risk zone of 100 m near landslides, and assessment of potential consequences.
            \item \textcolor{red}{Results}: calculated possible losses caused by the destruction of immobility and transport: costs for buildings demolition, restoration, roads rebuilding, debris transport, excavation and removal.
\end{enumerate}
\end{column}
\end{columns}
\end{minipage}
\end{frame}

\section{Geographic Settings}

\subsection{Study area: Ybbs valley}
\begin{frame}{Study area: Ybbs valley, Alpine Foreland, Austria}

\begin{minipage}[0.4\textheight]{\textwidth}
\begin{columns}[T]
\begin{column}{0.5\textwidth}
\vspace{2em}
\begin{figure}[H]
	\centering
		\includegraphics[width=4.5cm]{Fig_01.jpg}
\end{figure}
\begin{figure}[H]
	\centering
		\includegraphics[width=4.5cm]{Fig_02.jpg}
\end{figure}
\end{column}
\begin{column}{0.60\textwidth}
\begin{enumerate}[i]
\vspace{1em}
\footnotesize
        \item[--] The region of Ybbs valley is located in mountainous Alpine Foreland, Austria.
       \item[--] It covers a total area of 131.3 km2.
       \item[--] The environmental conditions of the the Ybbs valley are especially prone to landslides. 
       \item[--] Geology: deposits of gravel, sand, marl, conglomerate, calcareous sandstone, clay. 
       \item[--] Most of the region is used as pasture, agricultural land, and large forests.
       \item[--] Climatic humid to continental. Intensive rainfalls in spring and autumn months. 
       \item[--] The region is attractive for tourists due to its environmental and historical value.
\end{enumerate}
\vspace{0em}
\begin{figure}[H]
	\centering
		\includegraphics[width=5.5cm]{Fig_03.jpg}
\end{figure}
\end{column}
\end{columns}
\end{minipage}
\end{frame}

\section{Methods}
\subsection{Data} 
\begin{frame}\frametitle{Data}
\begin{minipage}[0.4\textheight]{\textwidth}
\begin{columns}[T]
\begin{column}{0.5\textwidth}
Data structure
\vspace{0em}
\begin{figure}[H]
	\centering
		\includegraphics[width=5.5cm]{Tab_01.jpg}
\end{figure}
\end{column}
\begin{column}{0.50\textwidth}
Data characteristics
\begin{enumerate}[i]
\vspace{1em}
\footnotesize
            \item \textcolor{red}{Sources}: Data used in this research were obtained from the repositories of Provincial Government of Lower Austria.
            \item \textcolor{red}{Content}: Data include vector thematic layers covering the area, aerial imagery, topographic vector layers and maps. 
            \item \textcolor{red}{Organization}: Data have been organized in a ArcCatalog database and sorted according to their types (Tab.1).
            \item \textcolor{red}{Compatibility}: Data have been stored in the ArcGIS project in Austrian National grid Austria Bundesmeldenetz (BMN) M34.
            \item \textcolor{red}{Software}: Arc GIS-based spatial analysis and calculations for estimation of the potential monetary losses caused by landslides.
\end{enumerate}
\end{column}
\end{columns}
\end{minipage}
\end{frame}


\section{Results}
\subsection{Mapping Results}
\begin{frame}\frametitle{Mapping Results}
\begin{minipage}[0.4\textheight]{\textwidth}
\begin{columns}[T]
\begin{column}{0.5\textwidth}
1988
\begin{figure}[H]
	\centering
		\includegraphics[width=5.0cm]{00.jpg}
\end{figure}
\end{column}
\begin{column}{0.5\textwidth}
2011
\begin{figure}[H]
	\centering
		\includegraphics[width=5.0cm]{00.jpg}
\end{figure}
\end{column}
\end{columns}
\end{minipage}
\end{frame}

\subsection{Conclusions}
\begin{frame}\frametitle{Conclusions}
\begin{itemize}
\footnotesize
	   \item The research focuses on the monetary estimation of the possible losses caused by landslide risks. Estimation of the economic damages is performed using existing simplified methodologies. 	   
            \item Calculations were based on real estate and market price of the elements at risk.            
            \item While assessing potential damage of landslides confusion arises due to these factors.
            \begin{itemize}
            \footnotesize
            	\item Temporal probability of the landslides occurrence is highly difficult to assess: it can only be estimated based on the reliable and obtainable data. This includes historical data continuously reporting the occurrence of the landslides.
		\item Difficulties arise by estimation of the indirect losses and partially damaged objects. The amount of the damages can be assessed based on elements vulnerability, which is very uncertain to estimate exactly. Thus, the vulnerability may differ depending on object location, individual characteristics and external factors.
		\item The term “landslide” include debris flows, shallow or rotational landslides. 
		\item Future studies may include different method approaches in the risk assessment to assess difference in magnitude and occurrence of landslides, risk perception and vulnerability assessment. 
	    \end{itemize}
            \item The data about landslide probability are gained from economic monitoring programmes.
            \item The elements at risk are defined based on spatial analysis and infrastructure inventory.
            \item The vulnerability estimation included basic census and approximate social data.
            \item Results confirm presence of the destructive processes of landslides.
            \item Research demonstrated successful Arc GIS based of the spatial data analysis and risk assessment
\end{itemize}
\end{frame}

\end{document}
%Changing the font size locally (from biggest to smallest):	
%\Huge
%\huge
%\LARGE
%\Large
%\large
%\normalsize (default)
%\small
%\footnotesize
%\scriptsize
%\tiny

\end{document}